\section{テストベンチの仕様：再現のために}
\label{sec:dishspec}

以下に、このテストベンチを組み直すのに必要なものをすべてまとめる。装置、ループ、入力の符号化、出力の読み出し、フィードバック、実験プロトコルである。各パラメータには出典を記す。「論文」は、値を論文\cite{Kagan2022}の本文から採ったことを示す。「選択」は、論文に記載がなく、再現するには自分で決めなければならないことを示す。既定値を提案し、それをどう確かめるかを述べる。

\subsection*{装置と培養}

\begin{center}
\footnotesize
\setlength{\tabcolsep}{4pt}
\renewcommand{\arraystretch}{1.15}
\begin{tabular}{>{\raggedright}p{4.0cm}>{\raggedright}p{6.9cm}>{\raggedright\arraybackslash}p{1.6cm}}
\toprule
パラメータ & 値 & 出典 \\
\midrule
チップ & MaxOne アレイ：$8$~mm$^2$ の区画に白金電極 $26\,000$ 本 & 論文 \\
記録 & 同時に最大 $1024$ チャネル、毎秒 $20\,000$ サンプル & 論文 \\
刺激 & 技術的には最大 $32$ 電極、実験では独立な $8$ 電極 & 論文 \\
細胞 & マウス胎児の大脳皮質。幹細胞由来のヒトニューロン（2通りの培養法）。対照として HEK293T 細胞（ニューロンではない） & 論文 \\
播種 & チップあたり約 $10^6$ 個 & 論文 \\
実験時の培養日数 & マウス：$\sim10$ 日から。ヒト：培養法により $14$--$24$ 日または $\sim80$ 日 & 論文 \\
\bottomrule
\end{tabular}
\end{center}

\subsection*{ループと時間}

ループは $10\ms$（$200$ サンプル）の窓ごとに動く。窓ごとに運動領域の電極のスパイクを数え、ゲームを更新し、次の刺激を出す。培養内のスパイクから応答刺激までの遅れは約 $5\ms$ である。プレイ中、各電極の信号は遮断周波数 $100$~Hz の2次ベッセル高域通過フィルタを通る。信号の絶対値は $1$~Hz の1次ベッセル低域通過フィルタで平滑化され、スパイクの閾値はこの平滑化した絶対値に比例する。背景の標準偏差の6倍という閾値を論文が挙げているのは、個別の活動測定についてであり、ゲームについてではない。刺激が培養の応答として数えられないよう、$75$~mV を超える大きなパルスが同時に $15$ 個より多く来たときは、すべての信号を消去する。以上はすべて論文による。

\subsection*{ボールの符号化}

\begin{center}
\footnotesize
\setlength{\tabcolsep}{4pt}
\renewcommand{\arraystretch}{1.15}
\begin{tabular}{>{\raggedright}p{3.4cm}>{\raggedright}p{7.5cm}>{\raggedright\arraybackslash}p{1.6cm}}
\toprule
パラメータ & 値 & 出典 \\
\midrule
感覚領域 & $8$ 本の電極。隣り合う電極が隣り合うボール位置を表すように配置 & 論文 \\
場所による符号 & 電極番号 $k=1,\dots,8$ がラケット中心に対するボールの位置を表す & 論文 \\
どの座標か & ボールの垂直位置。再現ではラケットに対する位置とし、$y-p\in[-\tfrac12,\tfrac12]$ で $k=\lceil 8(y-p+\tfrac12)\rceil$ & 論文。式は選択 \\
頻度による符号 & $4$ から $40$~Hz の刺激頻度がラケットまでの距離を運ぶ & 論文 \\
目盛りの向き & ボールが反対側の壁にあるとき $4$~Hz、ラケット側の壁に近づくにつれ線形に $40$~Hz まで：$f=4+36\,(1-x/L)$。$x$ はラケット側の壁までの距離、$L$ はコートの幅 & 論文 \\
パルス波形 & 短い矩形の二相性パルス。正の相が先、負の相が後 & 論文 \\
振幅 & $75$~mV。抑制された細胞を無理に発火させないように選ばれた & 論文 \\
各相の長さ & 記載なし。刺激システムの標準設定を使い、実験中は変えない & 選択 \\
\bottomrule
\end{tabular}
\end{center}

\subsection*{ラケットの読み出し}

論文では運動領域は二つあり、一方のスパイクはラケットを上へ、他方は下へ動かす。細胞密度と活動の違いを補正するため、領域の寄与には重みがつけられる\cite{Kagan2022}。正確な式は示されていない。再現には規則\eqref{eq:dishreadout}、すなわち $10\ms$ の窓ごとに $\Delta p=c\,(n_\uparrow-n_\downarrow)$ を使い、安静時の平均活動で領域をそろえる重みをつける（選択）。$n_\uparrow$ と $n_\downarrow$ を、ゲーム前に測ったその領域の窓あたり平均カウントで割るのである。係数 $c$ は、平均カウント1個分の差がラケットを1窓あたりコートの高さの $1\%$ だけ動かすように選ぶ（選択）。すると一方の領域が一定に優勢なら、ラケットは $1$~s でコートを横断する。

チップ上の領域の配置は論文本文に座標で与えられておらず、模式図があるだけである。再現には重ならない電極群が三つあれば足りる。中央の感覚電極8本と、その両側の記録電極群で、一方が「上」、他方が「下」である（選択）。

\subsection*{ゲーム}

コートの大きさ、ラケットの長さ、ボールの速さは論文に記載がない。ボールが一定の速さで飛び、壁で跳ね返るとあるだけである。再現には（選択）：コート $1\times1$、右の壁に長さ $0.2$ のラケット（モデル~\texttt{models/dishbrain\_model.py} と同じく $|y-p|<0.1$ でヒット）、ボールは $1$~s でコートを横切る。ラリーで一度も打ち返さずにミスしたものを「エース」、3回を超えて連続で打ち返したラリーを「長いラリー」と数える（論文）。

\subsection*{フィードバック}

\begin{center}
\footnotesize
\setlength{\tabcolsep}{4pt}
\renewcommand{\arraystretch}{1.15}
\begin{tabular}{>{\raggedright}p{2.6cm}>{\raggedright}p{8.3cm}>{\raggedright\arraybackslash}p{1.6cm}}
\toprule
イベント & 与えるもの & 出典 \\
\midrule
ヒット & 予測可能な刺激：感覚電極 $8$ 本すべてに同時に、$75$~mV、$100$~Hz、$100\ms$。その間、通常のボールの刺激は中断する & 論文 \\
ミス & 予測不能な刺激：$150$~mV、$5$~Hz、$4$~s、ランダムな時刻に、$8$ 本からランダムに選んだ電極へ。その後 $4$~s 刺激なし、ボールはランダムな方向にスタート & 論文 \\
ランダムさの法則 & 記載なし。再現では、パルスの時刻を平均頻度 $5$~Hz のポアソン過程とする & 選択 \\
\bottomrule
\end{tabular}
\end{center}

\subsection*{プロトコルと対照}

1セッションは $20$~min で、最初の $5$~min と最後の $15$~min を比べる（論文）。結果の指標は三つ。平均ラリー長、エースの割合、長いラリーの割合である。実験条件（論文）：
\begin{itemize}
\item \emph{刺激}：上に述べた完全なループ。
\item \emph{無音}：ヒットやミスの刺激の代わりに同じ時間まったく刺激せず、その後ボールはランダムな方向にスタートする。
\item \emph{フィードバックなし}：ミスの後、ボールは単に跳ね返って飛び続け、ボール位置の刺激は続く。
\item \emph{安静}：ゲームは進み、ラケットは活動に従って動くが、刺激はまったくない。
\item \emph{対照}：培地だけのチップ。HEK293T 細胞。「コンピュータの中の培養」、つまり細胞の代わりのシミュレーション。
\end{itemize}
主系列の標本数：マウス培養 $9$ 個（$101$ セッション）、ヒト培養 $11$ 個（$138$）、培地のみのチップ $6$ 枚（$80$）、安静の培養 $20$ 個（$42$）、シミュレーション $3$ 個（$38$）、計 $399$ セッション。フィードバックの系列では、1日 $3$ セッションを $3$ 日で $486$ セッション、条件は「刺激」「無音」「フィードバックなし」（$n=27$、$17$、$15$）。最初の $5$ 分から最後の $15$ 分への平均ラリー長の伸びは、生きた培養でだけ得られた。フィードバックの系列で時間とともに有意に伸びたのは「刺激」条件だけだった。セッションの終わりには「無音」条件も「フィードバックなし」を上回ったが、効果は小さく、3日間でセッション間の安定した学習はなかった\cite{Kagan2022}。

\subsection*{テストベンチの1サイクル}

$10\ms$ ごとに、テストベンチのコンピュータは次のことを行う。
\begin{enumerate}
\item 直前の窓における二つの運動領域のスパイク数 $n_\uparrow$、$n_\downarrow$ を読み、安静時のその領域の平均カウントで正規化する。
\item ラケットを $\Delta p=c\,(n_\uparrow-n_\downarrow)$ だけ動かし、コートの端で制限する。
\item ボールを1ステップ進め、上、下、左の壁で跳ね返す。
\item ボールが右の壁に達したら、$|y-p|<0.1$ ならヒットで、ラリーのカウントを1増やし、ヒットの刺激（$100\ms$）を与える。そうでなければミスで、ラリーを記録し、ミスの刺激（$4$~s）を与え、$4$~s 休止してボールを再スタートする。
\item そうでなければ、ボールの垂直位置から電極 $k$ を、ラケット側の壁までの距離から頻度 $f$ を選び、次の窓まで電極 $k$ に $75$~mV の二相性パルスを頻度 $f$ で与える。
\end{enumerate}
これだけあれば、公表されたものと互換なテストベンチを組み、本章の中心的な問いを確かめられる。培養を教えるのは予測可能性なのか、それともヒットとミスへの応答の単なる違いなのか。そのためには論文の三つの条件に、論文にない四つ目を加えるとよい。ミスの後に、予測不能な刺激と同じ長さとエネルギーの\emph{予測可能な}刺激を与える条件である。それでも培養が学習するなら、効いているのは予測可能性ではなく違いである。
